% Source-derived chapter 11: Characterization of  $p$-adic periods
% Original source: standard-conjectures-abelian-fourfolds
\chapter{Characterization of  $p$-adic periods}
\label{standard-conjectures-abelian-fourfolds-chapter-11}
\source{standard-conjectures-abelian-fourfolds}

\begin{remark}[Source section]
\label{standard-conjectures-abelian-fourfolds-chapter-11-source}
\source{standard-conjectures-abelian-fourfolds}
\source{standard-conjectures-abelian-fourfolds}
This chapter reproduces the corresponding section of the retrieved
LaTeX source, including its definitions, statements, examples, and
proofs. It has not yet been translated into Lean.
\notready
\end{remark}

% The source section title was: Characterization of  $p$-adic periods
\label{p-adic periods}
\source{standard-conjectures-abelian-fourfolds}
This section continues the comparison between the two quadratic spaces
$(V_{B,p},q_{B,p})$ and  $(V_{Z,p},q_{Z,p})$ which we initiated in the previous section. We keep notation from there, in particular we work under the Assumption \ref{ipos} which gives a well defined positive integer $i$ and we will make constant use of the field $F$  constructed in  Proposition \ref{propendo}.

The goal is to study the period matrix (relating  the two quadratic spaces $V_{B,p}$ and $V_{Z,p}$) given by Corollary \ref{cortsuji}. 
We compute  the filtration and the action of the Frobenius on these periods. Most importantly, we show that the elements of $B_\crys$ having the same behavior with respect to filtration and Frobenius are essentially unique (this fact characterizes the periods involved).

We   divide the analysis in two cases, depending  whether $F$ is unramified or not. 


\subsection*{Unramified case}

In this subsection we work under   Assumption  \ref{ipos} and we assume moreover that the field $F$  constructed in  Proposition \ref{propendo} is unramified over $\Q_p$.
This means that we have the inclusion $F \subset B_\crys$ and that the absolute Frobenius $\varphi$ of $B_\crys$  restricted to $F$ is the  non-trivial element of $\Gal(F/\Q_p)$.
\begin{definition}
\source{standard-conjectures-abelian-fourfolds}
\notready\label{defpi}
\source{standard-conjectures-abelian-fourfolds}
We define the $F$-vector subspace $P_i=P_i(F)$ of $B_\crys$ as the set of $\lambda_i \in B_\crys$ verifying the following properties:
\begin{enumerate}
\item $\varphi^2 (\lambda_i)= \lambda_i.$
\item $\lambda_i \in \Fil ^i B_\dR.$
\item $\varphi(\lambda_i) \in \Fil ^{-i} B_\dR.$
\end{enumerate}
\end{definition}


\begin{proposition}
\source{standard-conjectures-abelian-fourfolds}
\notready\label{lemmapi}
\source{standard-conjectures-abelian-fourfolds}
The $F$-vector space  $P_i(F)$ is of dimension one. Moreover a nonzero element $\lambda_i \in P_i$ verifies the following properties.
\begin{enumerate}
 \item\label{preciso1} $\lambda_i \cdot\varphi(\lambda_i) \in \Q^*_p.$
\source{standard-conjectures-abelian-fourfolds}
  \item\label{preciso4} $\lambda_i $ and $\varphi(\lambda_i)$ are invertible in $B_\crys.$
\source{standard-conjectures-abelian-fourfolds}
\item\label{preciso2} $\lambda \in \Fil ^i B_\dR-\Fil ^{i+1} B_\dR.$
\source{standard-conjectures-abelian-fourfolds}
\item\label{preciso3} $\varphi(\lambda) \in \Fil ^{-i} B_\dR-\Fil ^{-i+1} B_\dR.$
\source{standard-conjectures-abelian-fourfolds}
\end{enumerate}
\end{proposition}

\begin{proof}

By construction of $P_i$ we have   $ \lambda_i \cdot\varphi(\lambda_i) \in (B_\crys^{\varphi=\id} \cap \Fil^0B_\dR)$  and on the other hand we have  $(B_\crys^{\varphi=\id} \cap \Fil^0B_\dR)=\Q_p$, see Theorem \ref{javier}. What is missing to show (\ref{preciso1}) is that $\lambda_i \cdot\varphi(\lambda_i)\neq 0.$ It will follow from part (\ref{preciso4}).

  Consider the $\Q_p$-vector space $N=\Q_p^2$ with basis $e_1,e_2$ endowed with 
\begin{equation}\label{biagi}\varphi=\begin{pmatrix}
\source{standard-conjectures-abelian-fourfolds}
0 & 1 \\
1 & 0
\end{pmatrix}
\end{equation}
as Frobenius and with the following filtration
\[\Fil^{-i}N=N, \Fil^{-i+1}N =\Fil^i N=\Q_p\cdot e_2, \Fil^{i+1} N=0.  \]
Let us  check that $N$ is an admissible filtered $\varphi$-module.  The action of $\varphi$ on $\det N$, the maximal exterior power of $N$, is by multiplication by $-1$, which has valuation equal to $0$, and the only non trivial step of the induced filtration on $\det N$ is at $0$ as well. There are two lines  stable by $\varphi$ and none of those is $\Q_p\cdot e_2$ hence the only non trivial step of the induced filtration on these lines is at $-i$. On the other hand the action of $\varphi$ on these lines is by multiplication by $\pm 1$, which has valuation equal to $0$. As $i$ is  positive, the required inequality  $-i \leq 0$ holds true.
 
 As $N$ is an admissible $\varphi$-module, by Theorem \ref{tsuji} there exists a unique $\Gal_{\Q_p}$-representation $V$ such that $N=D(V)$. 
 Consider the identification $V \otimes B_\crys= N \otimes B_\crys$ given by  Theorem \ref{tsuji}. It allows to write an element of $V$ with respect to the basis $e_1,e_2$. Let $\begin{pmatrix} \alpha \\
\beta
 \end{pmatrix}$ be a fixed nonzero vector in $V$. As the Frobenius must act trivially on it we get from (\ref{biagi})  the relations
 \[ \alpha=  \varphi(\beta), \,
   \beta= \varphi(\alpha) \]
  which give the equality
  \begin{equation}  \alpha=  \varphi^2(\alpha).\end{equation}
  
    On the other hand, again by Theorem \ref{tsuji}, the vector $\begin{pmatrix} \alpha \\
\beta
 \end{pmatrix}$ lies in $ \Fil^0[M \otimes B_\dR]$. This means that it can be written as a linear combination $\gamma n + \delta e_2$ with $\gamma \in \Fil^{i}B_\dR, \delta \in \Fil^{-i} B_\dR$ and $n \in N$, which implies that $\alpha \in \Fil^{i}B_\dR$ and  $\beta \in \Fil^{-i}B_\dR$. 
 
 Altogether we have    $\alpha \in P_i$ and in particular  $P_i$ is a non empty space. Moreover, arguing as above, one sees that for all $t \in P_i$ the vector $\begin{pmatrix} t   \\
\varphi(t)
 \end{pmatrix}$ is in the vector space $V=[N \otimes B_\crys]^{\varphi=\id} \cap  \Fil^0[N \otimes B_\dR] $. This gives an isomorphism of $\Q_p$-vector spaces between $V$ and $P_i$. As $V$ has dimension $2$ the space $P_i$ is an $F$-vector space of dimension one.
 
 Now notice that  the one dimensional $\Q_p$-vector space of $p$-adic periods relating $\det N$ to $\det V$ is generated by an element in $B_\crys$ of the form $\alpha \beta (f-\varphi( f) )$, for any fixed $f \in F- \Q_p$. As $\det N\otimes B_\crys$ and $\det V\otimes B_\crys$ are free $B_\crys$-modules of rank one the element $\alpha \beta (f-\varphi( f) )$ must be invertible in $B_\crys$, hence $\alpha$ and  $\beta$ are invertible as well.
 
Finally, let us consider the  valuation $\nu$ from Theorem \ref{javier} associated to the filtration. We have showed the inequalities $\nu ( \alpha) \geq +i$ and  $\nu (  \beta)\geq -i$.
On the other hand, as the only non trivial step of the filtration on $\det M$ is  $0$, we deduce that $\nu ( \alpha \beta (f-\sigma( f) ))=0$ and hence $\nu ( \alpha \beta)=0$. This implies  that the above inequalities are actually equalities.
 \end{proof}


\begin{definition}
\source{standard-conjectures-abelian-fourfolds}
\notready\label{defqi}
\source{standard-conjectures-abelian-fourfolds}
Consider $V_{Z,+}$ and $V_{B,+}$ as constructed in Remark \ref{conventionF} and consider them as $F$-vector subspaces of $V_{Z,p}\otimes B_\crys$ via Corollary \ref{cortsuji}. 
We define the $F$-vector subspace $Q=Q(M)$ of $B_\crys$ as the set of $\lambda \in B_\crys$ such that
\[\lambda \cdot V_{B,+} \subset V_{Z,+}.\]
\end{definition}
\begin{remark}
\source{standard-conjectures-abelian-fourfolds}
\notready
The notation $Q(M)$ would like to suggest that the definition of $Q$ depends a priori on the motive $M$ from Theorem \ref{mainthm}. It will turn out that it actually only depends on $i$ and $F$.
\end{remark}
\begin{proposition}
\source{standard-conjectures-abelian-fourfolds}
\notready\label{piqi}
\source{standard-conjectures-abelian-fourfolds}
The $F$-vector subspaces $P_i,Q\subset B_\crys$ introduced in Definitions \ref{defpi} and \ref{defqi} coincide and they have dimension one.
\end{proposition}

\begin{proof}
 Following Remark \ref{conventionF} we have 
 \begin{equation} \label{eq:mirella0}V_{B,+} \cdot B_\crys =  V_{Z,+} \cdot B_\crys.\end{equation}
\source{standard-conjectures-abelian-fourfolds}
  As the  $F$-vector spaces $V_{B,+} $ and $ V_{Z,+}$ have dimension one   then $Q\subset B_\crys$ must have dimension one as well.
 
Fix a nonzero element $\lambda \in Q$. By definition of $Q$ there are nonzero vectors $v_B \in V_{B,+} $ and $v_Z \in V_{Z,+} $ such that
\begin{equation} \label{eq:mirella} \lambda \cdot v_{B} = v_{Z}  \end{equation}
\source{standard-conjectures-abelian-fourfolds}
which implies 
\begin{equation} \label{eq:mirella2} \varphi(\lambda) \cdot \varphi(v_{B}) = \varphi(v_{Z}).
\source{standard-conjectures-abelian-fourfolds}
\end{equation}
Now  $\varphi$ acts trivially on $V_{B,p}$ and,  because of the presence of algebraic classes, also on $V_{Z,p}$. Hence $\varphi$ acts as the non-trivial element of $\Gal(F/\Q_p)$ on $V_{B,p}\otimes F$ and  $V_{Z,p}\otimes F$. This implies that 
\begin{equation} \label{eq:mirella3} \varphi(v_{B})  \in V_{B,-} \hspace{0.5cm}  \textrm{and} \hspace{0.5cm} \varphi(v_{Z})  \in V_{Z,-} 
\source{standard-conjectures-abelian-fourfolds}
\end{equation}
as well as 
\begin{equation} \label{eq:mirella4} \varphi^2(v_B)=v_B \hspace{0.5cm}  \textrm{and} \hspace{0.5cm}   \varphi^2(v_Z)=v_Z.
\source{standard-conjectures-abelian-fourfolds}
\end{equation}
We claim that  the above relations (through Corollary \ref{cortsuji}) imply   that $\lambda$ verifies conditions (1),(2) and (3) from Definition \ref{defpi}. Indeed (\ref{eq:mirella}) gives (2), (\ref{eq:mirella2}) and  (\ref{eq:mirella3}) give (3) and finally  (\ref{eq:mirella4}) gives (1).

Hence, we have showed the  inclusion $Q \subset P_i$. For dimensional reasons (Proposition \ref{lemmapi}) it is actually an equality.
\end{proof}

\begin{prop}
\source{standard-conjectures-abelian-fourfolds}
\notready\label{warmbrinon}
\source{standard-conjectures-abelian-fourfolds}
Consider a nonzero element $\lambda_i \in P_i,$ and recall that $\lambda_i \cdot\varphi(\lambda_i) \in \Q^*_p$ (Proposition \ref{lemmapi}).  Then the quadratic forms $q_{B,p}$ and $q_{Z,p}$ are isomorphic if and only if
 \[\lambda_i \cdot\varphi(\lambda_i) \in N_{F/\Q_p}(F^*).\] 
\end{prop}
\begin{proof}
Let $v_B \in V_{B,+}$ and $v_Z \in V_{Z,+}$ as in the proof of Proposition \ref{piqi}.  By (\ref{eq:mirella4}) we have
\begin{equation} \label{eq:mirella9}v_B + \varphi(v_{B})  \in  V_{B,p} \hspace{0.5cm}  \textrm{and} \hspace{0.5cm}    v_Z + \varphi(v_{Z})  \in  V_{Z,p}. 
\source{standard-conjectures-abelian-fourfolds}
\end{equation}

Write  $\langle \cdot  ,  \cdot \rangle_{?}$ for the bilinear form associated to the quadratic form $q_?$. Then we have
\[q_{B,p}(v_B + \varphi(v_{B}) ) = 2 \langle v_B ,  \varphi(v_{B}) \rangle_{B,p}  \hspace{0.3cm}  \textrm{and}   \hspace{0.3cm} q_{Z,p}(v_Z + \varphi(v_{Z}) ) = 2 \langle v_Z ,  \varphi(v_{Z}) \rangle_{Z,p}  \]
because    $v_B,v_Z,\varphi(v_B)$ and $\varphi(v_Z)$ are isotropic vectors (see Remark \ref{conventionF} and (\ref{eq:mirella3})). 
Hence, using  relations  (\ref{eq:mirella}), (\ref{eq:mirella2}), together with (\ref{chissa}), we have
\[\lambda_i \cdot\varphi(\lambda_i) \cdot q_{B,p}(v_B + \varphi(v_{B}) ) =  q_{Z,p}(v_Z + \varphi(v_{Z}) ) . \]
We can conclude thanks to Corollary \ref{norms}.
\end{proof}

\subsection*{Ramified case}\label{ramcase}
\source{standard-conjectures-abelian-fourfolds}


In this subsection we work under the Assumption  \ref{ipos} and we assume moreover that the field $F$  constructed in  Proposition \ref{propendo} is ramified over $\Q_p$.
Define $\sigma$ to be the  non-trivial element of $\Gal(F/\Q_p)$ and $  B_{\crys,F} \subset B_\dR$ as the smallest subring containing $F$ and  $B_{\crys}$.
The inclusions   $  B_{\crys} \subset B_{\crys,F}$ and  $  F \subset B_{\crys,F}$ induce an identification
\[F\otimes_{\Q_p} B_{\crys}=B_{\crys,F}. \]
This allows  to extend $\sigma$ and $\varphi$ to two endomorphisms of $\Q_p$-algebras
\[   \sigma, \varphi : B_{\crys,F}    \longrightarrow B_{\crys,F} \]
which commute.

\begin{definition}
\source{standard-conjectures-abelian-fourfolds}
\notready\label{defpiram}
\source{standard-conjectures-abelian-fourfolds}
We define the $F$-vector subspace $P_i=P_i(F)$ of $B_{\crys,F}$ as the set of $\lambda \in B_{\crys,F}$ verifying the following properties.
\begin{enumerate}

\item $\lambda \in \Fil ^i B_\dR.$
\item $\sigma(\lambda) \in \Fil ^{-i} B_\dR.$
\item $\varphi (\lambda)= \lambda.$
\end{enumerate}
\end{definition}

\begin{proposition}
\source{standard-conjectures-abelian-fourfolds}
\notready\label{lemmapiram} The $F$-vector space  $P_i(F)$ is of dimension one. 
\source{standard-conjectures-abelian-fourfolds}
Moreover a nonzero element $\lambda_i \in P_i$  verifies the following properties.
\begin{enumerate}
 \item  $\lambda_i \cdot\sigma(\lambda_i) \in \Q^*_p.$
  \item  $\lambda_i$ and $\sigma(\lambda_i)$ are invertible in $B_{\crys,F}.$
\item  $\lambda_i \in \Fil ^i B_\dR-\Fil ^{i+1} B_\dR.$
\item  $\sigma(\lambda_i) \in \Fil ^{-i} B_\dR-\Fil ^{-i+1} B_\dR.$
\end{enumerate}
\end{proposition}


\begin{proof}
By construction of $P_i$ we have that $\lambda_i \cdot\varphi(\lambda_i) $ is invariant under $\sigma$ hence it is in $B_\crys$. Moreover it verifies  $ \lambda_i \cdot\varphi(\lambda_i) \in (B_\crys^{\varphi=\id} \cap \Fil^0B_\dR)$  and on the other hand we have  $(B_\crys^{\varphi=\id} \cap \Fil^0B_\dR)=\Q_p$, see Theorem \ref{javier}. What is missing to show (\ref{preciso1}) is that $\lambda_i \cdot\varphi(\lambda_i)\neq 0.$ It will follow from part (\ref{preciso4}).

 Consider the $\Q_p$-vector space $N=\Q_p^2$ endowed with the identity as Frobenius. Let $L$ be a line in $N\otimes F$ which is not stable under $\sigma$.
Define the following filtration
\[\Fil^{-i}N\otimes F=N\otimes F, \Fil^{-i+1}N\otimes F =\Fil^i N \otimes F=\sigma(L), \Fil^{i+1} N=0.  \]
Let us  check that $N$ is an admissible filtered $\varphi$-module.  The action of $\varphi$ on $\det N$, the maximal exterior power of $N$, is by multiplication by $+1$, which has valuation equal to $0$, and the only non trivial step of the induced filtration on $\det N$ is at $0$ as well. Consider a line $L'$ in $N$. The action of $\varphi$ on it is again by multiplication by $+ 1$, which has valuation equal to $0$, and the only non trivial step of the induced filtration on it is at $-i$ because $L'\neq \sigma(L)$. As $i$ is  positive, the required inequality  $-i \leq 0$ holds true.
 
 As $N$ is an admissible $\varphi$-module, by Theorem \ref{tsuji} there exists a unique $\Gal_{\Q_p}$-representation $V$ such that $N=D(V)$. Consider the action of $f\in F$ on $N\otimes F$ given by multiplication by $f$ on $L$ and by multiplication by $\sigma(f)$ on $ \sigma(L)$. It descends over $\Q_p$ and gives an action of $F$ on $N$. Moreover this action respects the filtration. By  Theorem \ref{tsuji}, the field $F$ acts on $V$ as well and  the identification $V \otimes B_\crys= N \otimes B_\crys$ is $F$-equivariant. In particular there are eigenvectors 
 $m\in L$ and  $v \in V \otimes F$ and a scalar $\alpha \in B_{\crys,F}$ such that 
  \[
\alpha \cdot m=   v
 \]
 and hence $ \sigma(\alpha) \cdot \sigma(m)=   \sigma(v)$.
 
As the Frobenius acts trivially on $V$ and $N$ we deduce that it acts as the identity on  $\alpha$ and $\sigma(\alpha)$ as well. 
Moreover,  the inclusion $V \subset \Fil^0 N\otimes B_\dR$ implies $\alpha \in \Fil ^i B_\dR-\Fil ^{i+1} B_\dR$ and $\sigma(\alpha) \in \Fil ^{-i} B_\dR-\Fil ^{-i+1} B_\dR.$ 

 Altogether we have    $\alpha \in P_i$ and in particular  $P_i$ is a non empty space. Moreover, arguing as above, one sees that for all $t \in P_i$ the linear combination $t \cdot m + \sigma(t) \cdot \sigma (m)$  is in the vector space $V=[N \otimes B_\crys]^{\varphi=\id} \cap  \Fil^0[N \otimes B_\dR] $. This gives an isomorphism of $\Q_p$-vector spaces between $V$ and $P_i$. As $V$ has dimension $2$ the space $P_i$ is an $F$-vector space of dimension one.
 
 Finally, as $v$ and $m$ generates the same free  $B_{\crys,F}$-module of rank one, the scalar $\alpha$ must be invertible in $B_{\crys,F}$, and the same argument applies for $ \sigma(\alpha)$.
\end{proof}

\begin{definition}
\source{standard-conjectures-abelian-fourfolds}
\notready\label{defqiram}
\source{standard-conjectures-abelian-fourfolds}
Consider $V_{Z,+}$ and $V_{B,+}$ as constructed in Remark \ref{conventionF} and consider them as $F$-vector subspaces of $V_{Z,p}\otimes B_{\crys,F}$ via Corollary \ref{cortsuji}. 
We define the $F$-vector subspace $Q=Q(M)$ of $B_{\crys,F}$ as the set of $\lambda \in B_{\crys,F}$ such that
\[\lambda \cdot V_{B,+} \subset V_{Z,+}.\]
\end{definition}
\begin{proposition}
\source{standard-conjectures-abelian-fourfolds}
\notready\label{piqiram}
\source{standard-conjectures-abelian-fourfolds}
The $F$-vector subspaces $P_i,Q\subset B_{\crys,F}$ introduced in Definitions \ref{defpiram} and \ref{defqiram} coincide and they have dimension one.
\end{proposition}
\begin{proof}
Analogous to the proof of Proposition \ref{piqi}.
\end{proof}


\begin{prop}
\source{standard-conjectures-abelian-fourfolds}
\notready\label{warmbrinonram}
\source{standard-conjectures-abelian-fourfolds}
Consider a nonzero element $\lambda_i \in P_i,$ then \[\lambda_i \cdot\sigma(\lambda_i) \in \Q^*_p.\] Moreover the quadratic forms $q_{B,p}$ and $q_{Z,p}$ are isomorphic if and only if
 \[\lambda_i \cdot\sigma(\lambda_i) \in N_{F/\Q_p}(F^*).\] 
\end{prop}
\begin{proof}
Analogous to the proof of Proposition \ref{warmbrinon}.
\end{proof}
\begin{remark}
\source{standard-conjectures-abelian-fourfolds}
\notready  
For any $i$ and any $F$ the periods $P_i(F)$ appear in the realization of a motive $M_{i,F}$. Indeed, for $i=1$ this comes from
Proposition \ref{propexemple}, see also Remark \ref{bellaciao}(\ref{unamattina}). For a different $i$ one can consider $M_{i,F}=M_{1,F}^{\otimes i}$, where the tensor product is done in the $F$-linear category of motives endowed with an $F$-action. 

Define $P_0(F)=F$ and $P_*(F)=\bigoplus_{i\geq 0} P_i(F)$. By the very definition of $P_i(F)$, the set $P_*(F)$ is a graded $F$-algebra. Moreover Propositions \ref{lemmapi} and \ref{lemmapiram} imply that $P_*(F)$  is non canonically isomorphic to the polynomial algebra $F[x]$.

\end{remark}
